\Needspace{4\baselineskip}
\begin{Example}
\label{ex:integral-kernel-operator}
用下式定义\textbf{积分算子} $K:L^2(\Omega)\to L^2(\Omega)$:
\begin{equation}\label{eq:integral-kernel-operator-definition}
Ku:=\int_{\Omega}k(\bm{x},\bm{y})u(\bm{y})\,\mathrm{d}\bm{y}.
\end{equation}
这里 $k:\Omega\times\Omega\to\mathbb{C}$ 称为算子的\textbf{核}.

可以对核作若干假设, 以保证积分算子例如有界. 现在给出两个引理, 其所基于的假设在应用中将很重要.
\end{Example}

\begin{Lemma}
\label{lem:integral-kernel-row-column-boundedness}
设 $\Omega\subset\mathbb{R}^n$ 为有界区域, 并假设核 $k$ 满足
\begin{equation}\label{eq:integral-kernel-uniform-row-bound}
k_1:=\sup_{\bm{x}\in\Omega}\int_{\Omega}|k(\bm{x},\bm{y})|\,\mathrm{d}\bm{y}<\infty,
\end{equation}

\begin{equation}\label{eq:integral-kernel-uniform-column-bound}
k_2:=\sup_{\bm{y}\in\Omega}\int_{\Omega}|k(\bm{x},\bm{y})|\,\mathrm{d}\bm{x}<\infty
\end{equation}
那么由 \eqref{eq:integral-kernel-operator-definition} 定义的算子 $K:L^2(\Omega)\to L^2(\Omega)$ 有界.
\end{Lemma}

\begingroup
\postdisplaypenalty=10000
\renewcommand{\endProof}{\par\nobreak\hfill$\blacksquare$\par\vspace{0.1cm}}
\begin{Proof}
只需利用 Hölder 不等式得到\[\begin{aligned}
\|Ku\|_{L^2(\Omega)}^2&=\int_{\Omega}\left(\int_{\Omega}k(\bm{x},\bm{y})u(\bm{y})\,\mathrm{d}\bm{y}\right)^2\,\mathrm{d}\bm{x}\\
&\leq\int_{\Omega}\left(\int_{\Omega}\sqrt{|k(\bm{x},\bm{y})|}\sqrt{|k(\bm{x},\bm{y})|}u(\bm{y})\,\mathrm{d}\bm{y}\right)^2\,\mathrm{d}\bm{x}\\
&\leq\int_{\Omega}\left\{\int_{\Omega}|k(\bm{x},\bm{y})|\,\mathrm{d}\bm{y}\right\}\left\{\int_{\Omega}|k(\bm{x},\bm{y})|u(\bm{y})^2\,\mathrm{d}\bm{y}\right\}\,\mathrm{d}\bm{x}\\
&\leq k_1\int_{\Omega}\int_{\Omega}|k(\bm{x},\bm{y})|u(\bm{y})^2\,\mathrm{d}\bm{y}\,\mathrm{d}\bm{x}\\
&=k_1\int_{\Omega}u(\bm{y})^2\int_{\Omega}|k(\bm{x},\bm{y})|\,\mathrm{d}\bm{x}\,\mathrm{d}\bm{y}\\
&\leq k_1k_2\|u\|_{L^2(\Omega)}^2.
\end{aligned}\]证明完毕.
\end{Proof}
\endgroup

\begin{Lemma}
\label{lem:hilbert-schmidt-kernel-boundedness}
设 $\Omega\subset\mathbb{R}^n$ 为区域. 设 $k$ 为 \textbf{Hilbert-Schmidt} 的, 即假设
\begin{equation}\label{eq:hilbert-schmidt-kernel-square-integrability}
\int_{\Omega}\int_{\Omega}|k(\bm{x},\bm{y})|^2\,\mathrm{d}\bm{x}\,\mathrm{d}\bm{y}=C<\infty.
\end{equation}
那么由 \eqref{eq:integral-kernel-operator-definition} 定义的算子 $K:L^2(\Omega)\to L^2(\Omega)$ 有界.
\end{Lemma}

\begingroup
\postdisplaypenalty=10000
\renewcommand{\endProof}{\par\nobreak\hfill$\blacksquare$\par\vspace{0.1cm}}
\begin{Proof}
再次利用 Hölder 不等式得到\[\begin{aligned}
\|Ku\|_{L^2(\Omega)}^2&=\int_{\Omega}\left(\int_{\Omega}k(\bm{x},\bm{y})u(\bm{y})\,\mathrm{d}\bm{y}\right)^2\,\mathrm{d}\bm{x}\\
&\leq\int_{\Omega}\left\{\int_{\Omega}|k(\bm{x},\bm{y})|^2\,\mathrm{d}\bm{y}\right\}\left\{\int_{\Omega}u(\bm{y})^2\,\mathrm{d}\bm{y}\right\}\,\mathrm{d}\bm{x}\\
&=\left[\int_{\Omega}\int_{\Omega}|k(\bm{x},\bm{y})|^2\,\mathrm{d}\bm{y}\,\mathrm{d}\bm{x}\right]\int_{\Omega}u(\bm{y})^2\,\mathrm{d}\bm{y}\\
&=C\|u\|_{L^2(\Omega)}^2.
\end{aligned}\]证明完毕.
\end{Proof}
\endgroup

在 Section~\ref{sec:greens-functions} 中已经遇到核为 Green 函数的积分算子. 将在 Section~\ref{sec:integral-operators} 中更详细地研究积分算子.
